\section[Conjugate quotients and minimal memory]{Conjugate quotients and
minimal memory of the finite atlas}
\label{sec:cocientes-conjugados-rev11}
\index{quotient!conjugate}
\index{minimal memory!}
\index{internal selector}
\index{dodecaphonic word}

The double projection requires the transported state to be distinguished from each of its images. At the finite level of the atlas, this distinction admits an exact formulation: the same dodecaphasic word has several quotients, each of which preserves a given family of invariants and contracts the remainder. The configuration historically called \emph{diamond} is formulated here as a system of conjugate quotients over a common support. Its mathematical content consists in identifying the fibers of those quotients and determining the minimal memory required to separate the states that a projection has identified.

\subsection{Cellular domain, routes, and enriched internal state}

Let
\[
 \mathcal C_9=\{1,\ldots,9\}^2
\]
the finite APP chart. For a cell \(x=(r,c)\), the atlas fixes the internal
cellular type, the level
\[
 G(r,c)=
 \begin{cases}
 0,&(r,c)=(5,5),\\
 \max\{1,|r-5|,|c-5|\},&\text{otherwise},
 \end{cases}
\]
and, in the corresponding sector, the colour class \((r-c)\bmod 3\). The
published multiplicative evaluation is
\[
 a_{\times}(r,c)=\dr(rc).
\]
These readers are computed from the oriented cell; a later physical denomination
does not intervene in their definition.

The route catalog associates with \(x\) the coordinate
\[
 \mathcal R(x)=
 \bigl(n_A,s_C,k_{\Delta_4},\nu_{120},\nu_{270},\tau_{12}\bigr)
 \in\mathscr R_{56},
\]
where \(\tau_{12}\in\{-1,0,1\}^{12}\) is the dodecaphasic word. The
exhaustive comparison of the atlas of \(81\) cells with the expanded CT--108
catalogue establishes a unique surjective map
\[
 \mathcal R:\mathcal C_9\longrightarrow\mathscr R_{56},
 \qquad |\operatorname{im}\mathcal R|=56,
\]
whose fibers have the distribution
\begin{equation}
 41\cdot1+10\cdot2+2\cdot3+1\cdot4+2\cdot5=81.
 \label{eq:fibras-r56-rev11}
\end{equation}

For \(\tau=(\tau_0,\ldots,\tau_{11})\) and
\(d\in\{3,4,6\}\), let \(\mathcal O_d\) be the orbits of
\(j\mapsto j+d\) in \(\mathbb Z/12\mathbb Z\). We define
\[
 \nu_d(\tau)=
 \sum_{O\in\mathcal O_d}
 \operatorname{sgn}\!\left(\sum_{j\in O}\tau_j\right),
 \qquad
 (\nu_{90},\nu_{120},\nu_{180})=(\nu_3,\nu_4,\nu_6),
\]
and the cyclic correlations
\[
 Q_k(\tau)=\sum_{j=0}^{11}\tau_j\tau_{j+k},
 \qquad k=1,\ldots,6,
\]
with indices taken modulo twelve. They complete the reader.
\[
 z_0(\tau)=\#\{j:\tau_j=0\},\qquad
 w(\tau)=\sum_j|\tau_j|,\qquad
 L(\tau)=\sum_j\tau_j.
\]
The notation \(z_0\) separates the number of null entries from the distinct invariant
\(R_3=Q_3/2\).

\HMTClaim{HMT-R-ATLAS-MINIMAL-INTERNAL-SELECTOR}
\begin{theorem}[Minimal internal selector]
\label{thm:selector-interno-minimo-rev11}
Having fixed the origin of the chart and its lexicographic order, let
\[
 \kappa(x)=
 \#\{y\in\mathcal R^{-1}(\mathcal R(x)):y<_{\rm lex}x\}
 \in M_5:=\{0,1,2,3,4\}.
\]
The map
\begin{align}
 \mathcal S_{\rm int}:\mathcal C_9&\longrightarrow
 \mathscr E_{\rm int}:=\mathscr R_{56}\times M_5\times\mathbb Z^9,
 \notag\\
 x&\longmapsto
 \bigl(\mathcal R(x),\kappa(x),\nu_{90}(\tau_x),
 \nu_{180}(\tau_x),Q_1(\tau_x),\ldots,Q_6(\tau_x),z_0(\tau_x)\bigr)
 \label{eq:selector-interno-rev11}
\end{align}
is injective. Among the injective lifts of the form
\((\mathcal R,m)\), its memory alphabet has minimal cardinality.
\end{theorem}

\begin{proof}
The distribution \eqref{eq:fibras-r56-rev11} contains fibers of cardinality five. By the pigeonhole principle, every map \(m:\mathcal C_9\to M\) for which \((\mathcal R,m)\) is injective satisfies \(|M|\geq5\). Within each fiber, \(\kappa\) enumerates its elements exactly according to the fixed order and therefore separates two cells sharing a route. Cells with distinct routes are already separated by \(\mathcal R\). Consequently, \((\mathcal R,\kappa)\), and a fortiori \(\mathcal S_{\rm int}\), is injective and attains the lower bound.
\end{proof}

Projection onto the first coordinate yields the commutative diagram.
\[
\begin{CD}
 \mathcal C_9 @>{\mathcal S_{\rm int}}>> \mathscr E_{\rm int}\\
 @V{\mathcal R}VV @VV{\operatorname{pr}_{\mathscr R}}V\\
 \mathscr R_{56} @>{\operatorname{id}}>> \mathscr R_{56}.
\end{CD}
\]
Thus, \(\kappa\) measures the minimum cellular memory forgotten by
\(\mathcal R\). Its canonicity is relative to the origin and ordering of the chart:
if these are regarded as a gauge, the intrinsic object is the ordered fiber
or its orbit. The index \(\kappa\) and the antisheet memory
\(A_6=(e_0-e_6,\ldots,e_5-e_{11})\) answer to losses of information of
the same categorical nature, but fulfill different functions:
\(\kappa\) separates cells of a fiber of \(\mathcal R\), whereas \(A_6\)
participates in reconstructing the dodecaphasic register.

\subsection{Thick types and natural multisections}

The ratio
\[
 q:\mathcal C_9\longrightarrow\mathscr Y_{13},
 \qquad
 q(x)=\bigl(\operatorname{familia}(x),G(x)\bigr),
\]
has thirteen values. The cardinalities of its fibers are
\[
 1,8,16;\qquad 2,4,6,8;\qquad 2,4,6,8;\qquad 4,12.
\]
The output determined by \(y\in\mathscr Y_{13}\) is the finite multisection
\begin{equation}
 \mathfrak S(y)=
 \{\mathcal S_{\rm int}(x):q(x)=y\}
 \in\operatorname{Fin}^{+}(\mathscr E_{\rm int}).
 \label{eq:multiseccion-natural-rev11}
\end{equation}

Let \(\rho(r,c)=(10-r,10-c)\). The thirteen fibers are stable under
\(\rho\), but only \(q^{-1}(q(5,5))\) contains a fixed point. Therefore, in
the remaining twelve fibers, a pointwise \(\rho\)-equivariant section does not
exist: the multisection \eqref{eq:multiseccion-natural-rev11} is the natural output
of the coarse type. Choosing a representative requires additional
orientation or representation data and is not obtained from the name
subsequently assigned to the type.

\subsection{Dodecaphasic conjugation and signature quotients}

Let \(\mathcal T_{12}=\{-1,0,1\}^{12}\) and define the involution
\[
 C_M:\mathcal T_{12}\longrightarrow\mathcal T_{12},
 \qquad C_M(\tau)=-\operatorname{rev}(\tau).
\]
We group the linear and quadratic readers together.
\[
 \mathcal L(\tau)=
 (\nu_{90},\nu_{120},\nu_{180},L)\in\mathbb Z^4,
\]
\[
 \mathcal Q(\tau)=
 (Q_1,\ldots,Q_6,z_0,w)\in\mathbb Z^8.
\]
Conjugation satisfies the identities
\begin{equation}
 \mathcal L(C_M\tau)=-\mathcal L(\tau),
 \qquad
 \mathcal Q(C_M\tau)=\mathcal Q(\tau).
 \label{eq:paridad-conjugacion-rev11}
\end{equation}
For any reader \(F:\mathcal T_{12}\to I_F\), write
\[
 \tau\sim_F\tau'\quad\Longleftrightarrow\quad F(\tau)=F(\tau'),
 \qquad
 \mathfrak Q_F=\mathcal T_{12}/\!\sim_F.
\]
When \(F\) is equivariant with respect to \(C_M\), the involution descends to
\(\overline C_{M,F}\), and one obtains
\[
\begin{CD}
 \mathcal T_{12} @>{C_M}>> \mathcal T_{12}\\
 @V{q_F}VV @VV{q_F}V\\
 \mathfrak Q_F @>{\overline C_{M,F}}>> \mathfrak Q_F.
\end{CD}
\]
In particular, \(\overline C_{M,\mathcal Q}\) is the identity and \(\overline C_{M,\mathcal L}\) changes the
sign of the published signature. This parity difference is the conjugate
structure of the two quotients; they are not two evaluations of distinct
initial states.

The atlas contains seven word forms of dodecaphonic with multiplicities.
\begin{equation}
 54,2,2,2,7,7,7.
 \label{eq:multiplicidades-formas-rev11}
\end{equation}
One form is fixed by \(C_M\); the other six constitute three
conjugate pairs. In each pair, the cellular multiplicities are \(2\) and \(7\).
Consequently, no bijection of \(\mathcal C_9\) can lift
\(C_M\) while preserving the form fibers. The inversion \(\rho\) intertwines the
words in \(49\) cells and leaves \(32\) exceptions. The exact anti-section
at this level is therefore the correspondence
\[
 x\longmapsto
 \{y\in\mathcal C_9:\tau_y=C_M(\tau_x)\}.
\]
The numerical equality \(54+9+9+9=81\) describes the grouping of one fixed
form and three pairs of multiplicity \(2+7\); it does not constitute a
decomposition into bijective orbits of the cells.

\subsection{Exact Configuration of Quotients}

Let us define the partial readers
\[
 \Sigma_7(\tau)=
 (\nu_{90},\nu_{120},\nu_{180},Q_3,Q_4,Q_6,z_0)
 \in\mathbb Z^7,
\]
\[
 \ell_3(\tau)=(\nu_{90},\nu_{120},\nu_{180})\in\mathbb Z^3,
 \qquad
 p_4(\tau)=(Q_3,Q_4,Q_6,z_0)\in\mathbb Z^4.
\]
Each reader induces its quotient \(\mathfrak Q_{\Sigma_7}\),
\(\mathfrak Q_{\ell_3}\), or \(\mathfrak Q_{p_4}\). The following three collisions
locate precisely what information each projection retains.

\HMTClaim{HMT-R-ATLAS-CONJUGATE-QUOTIENTS}
\begin{theorem}[Configuration of Conjugate Quotients]
\label{thm:configuracion-cocientes-conjugados-rev11}
In the finite catalog the following assertions hold simultaneously.
\begin{enumerate}[label=\textup{(\roman*)},leftmargin=2.2em]
 \item The boundary words \(H_\alpha\) and \(H_{\rm APP}\) are distinct
 and belong to distinct dihedral orbits, but
 \[
  \Sigma_7(H_\alpha)=\Sigma_7(H_{\rm APP})
  =(-3,-4,-5,3,2,2,6).
 \]
 The omitted coordinate of the nearest neighbor separates them:
 \[
  Q_1(H_\alpha)=4,
  \qquad Q_1(H_{\rm APP})=3.
 \]

 \item The obstructing support \(P_\pi\) and the template historically
 designated by \(\tau_{\rm lep}\) satisfy
 \[
  \ell_3(P_\pi)=\ell_3(\tau_{\rm lep})=(-3,-4,-6),
 \]
 while
 \[
  p_4(P_\pi)=(7,6,6,3),
  \qquad p_4(\tau_{\rm lep})=(4,3,2,5).
 \]

 \item The templates nominally inherited as electronic and protonic
 publish the same word,
 \[
  \tau_e^{\rm hist}=\tau_p^{\rm hist}=0^{12}.
 \]
 Every invariant that factors exclusively through
 \(\mathcal T_{12}\) necessarily takes the same value on both.
\end{enumerate}
Therefore, the complete configuration is a fan of quotients of the same
enriched state,
\begin{equation}
 \mathscr E_{\rm int}
 \longrightarrow
 \left\{
 \begin{array}{c}
  \mathfrak Q_{\mathcal L},\\[1mm]
  \mathfrak Q_{\mathcal Q},\\[1mm]
  \mathfrak Q_{\Sigma_7},\\[1mm]
  \mathfrak Q_{\ell_3},\\[1mm]
  \mathfrak Q_{p_4},\\[1mm]
  \mathcal T_{12},
 \end{array}
 \right.
 \label{eq:abanico-cocientes-rev11}
\end{equation}
and not a bijective correspondence between mathematical constants and physical classes.
\end{theorem}

\begin{proof}
The equalities in (i)--(iii) result from evaluating the readers defined above on the published words. In (i), the inequality of \(Q_1\) proves that the coincidence at \(\mathfrak Q_{\Sigma_7}\) arises from the suppressed coordinate and not from an identity of words. In (ii), the linear signature coincides and the even vector differs, so the identification occurs only in \(\mathfrak Q_{\ell_3}\). In (iii), equality of the words implies \(F(\tau_e^{\rm hist})=F(\tau_p^{\rm hist})\) for every map \(F\) whose effective domain is exclusively \(\mathcal T_{12}\). Thus each equality determines a specific fiber, and each separation identifies the memory that must be restored. This proves interpretation \eqref{eq:abanico-cocientes-rev11}.
\end{proof}

The theorem provides a finite model of the general principle of double
projection: an image may preserve an exact signature while also
identifying different states. The pair formed by the published signature and
the omitted memory refines the quotient; recovery of the state requires
reincorporating the cell, orientation, band section, or the pertinent coordinate
in each case.

\subsection{Accuracy of Conjugate Quotients and Obstructions to Memoryless Reconstruction}

The construction begins with the finite atlas of \(81\) cells, its \(56\)
route families, the seven section forms, and the involution \(C_M\), all
defined on the dodecaphase calendar of \(108\) steps. A state
is the quadruple comprising cell, route, orientation, and band section; this
quadruple precedes every physical designation and fixes the domain of the
conjugate quotients.

On that domain, the index \(\kappa\) and the minimality of five states follow from the origin and order of the chart. The equalities in \cref{thm:configuracion-cocientes-conjugados-rev11} are exact equalities of signatures: in each, the preserved reader and the memory coordinate forgotten by the quotient are specified. The exhaustive traversal of the \(81\) cells independently reproduces the fibers, the values of \(\nu_d\), the quotients \(Q_k\), and the diamond configuration; that finite enumeration complements the algebraic proof and permits its reproduction without altering its premises.

The conclusion is proved for these explicit primitives. Two extensions require additional maps: reconstructing the atlas from the minimal APP kernel and constructing, for the noncentral multisections, a fine physical section that preserves orientation and representation without consulting a target mass. Neither extension intervenes in the finite equalities already established.

The construction is incompatible with any of the following
variations: a cardinality different from \(81\), an image different from \(56\)
families, a fiber distribution different from
\eqref{eq:fibras-r56-rev11}, a failure in the reproduction of
\(\nu_d\) and \(Q_k\), the injectivity of
\(\mathcal S_{\rm int}\), the three equalities of
\cref{thm:configuracion-cocientes-conjugados-rev11}, the count
\(49+32\) of the partial intertwining of \(\rho\), or the functorial closure of
the cobordisms of the supplementary graph. A physical interpretation further requires
a representation map preserving these identities and a
comparison independent of its construction.
